\documentclass[11pt]{article}
\usepackage{amsmath,amssymb,amsthm}
\usepackage{algorithm}
\usepackage{algpseudocode}
\usepackage{geometry}
\geometry{margin=1in}
\newtheorem{theorem}{Theorem}
\newtheorem{lemma}{Lemma}
\newtheorem{assumption}{Assumption}
\newcommand{\E}{\mathbb{E}}
\newcommand{\R}{\mathbb{R}}
\newcommand{\norm}[1]{\left\|#1\right\|}
\newcommand{\ip}[2]{\left\langle #1,#2\right\rangle}
\begin{document}

\title{Convergence Analysis of a Noisy Gradient Descent Scheme \\ (SGLD‑style) }
\author{}
\date{}
\maketitle

%%%%%%%%%%%%%%%%%%%%%%%%%%%%%%%%%%%%%%%%%%%%%%%%%%%%%%%%%%%%
\section*{Algorithm Name}
\textbf{Noisy Gradient Descent with Decreasing Temperature (NGD‑DT)}  

The method updates the iterate $x_k\in\R^d$ by a deterministic gradient step
plus an additive Gaussian perturbation whose variance shrinks with the
iteration index.

%%%%%%%%%%%%%%%%%%%%%%%%%%%%%%%%%%%%%%%%%%%%%%%%%%%%%%%%%%%%
\section*{Mathematical Formulation}
Let $f:\R^d\rightarrow\R$ be the objective function.  
For a prescribed step‑size schedule $\{\alpha_k\}_{k\ge 0}$ and a
noise‑variance schedule $\{\sigma_k^2\}_{k\ge 0}$ we define
\[
\boxed{\;
x_{k+1}=x_k-\alpha_k \nabla f(x_k)+\xi_k,\qquad 
\xi_k\sim\mathcal{N}\!\bigl(0,\sigma_k^2 I_d\bigr)
\;}
\tag{1}
\]
All random variables are defined on a common probability space
$(\Omega,\mathcal{F},\mathbb{P})$.  The sequences $\{\alpha_k\}$ and
$\{\sigma_k^2\}$ are deterministic and non‑increasing.

%%%%%%%%%%%%%%%%%%%%%%%%%%%%%%%%%%%%%%%%%%%%%%%%%%%%%%%%%%%%
\section*{Pseudocode}
\begin{algorithm}[H]
\caption{NGD‑DT}
\begin{algorithmic}[1]
\State \textbf{Input:} initial point $x_0\in\R^d$, step‑size schedule
$\{\alpha_k\}$, noise schedule $\{\sigma_k^2\}$, maximal iteration $K$.
\For{$k=0,1,\dots,K-1$}
    \State Compute the gradient $g_k\gets \nabla f(x_k)$.
    \State Sample $\xi_k\sim\mathcal{N}(0,\sigma_k^2 I_d)$.
    \State Update $x_{k+1}\gets x_k-\alpha_k g_k+\xi_k$.
\EndFor
\State \textbf{Output:} $\{x_k\}_{k=0}^{K}$ (or the average iterate
$\bar x_K:=\frac1K\sum_{k=0}^{K-1}x_k$).
\end{algorithmic}
\end{algorithm}

%%%%%%%%%%%%%%%%%%%%%%%%%%%%%%%%%%%%%%%%%%%%%%%%%%%%%%%%%%%%
\section*{Dry Run Example}
Consider the one‑dimensional quadratic
\[
f(w)=(w-3)^2,\qquad \nabla f(w)=2(w-3).
\]
Set $w_0=0$, constant step size $\alpha=0.1$ and a deterministic
noise $\xi_k=0$ (for illustration we ignore stochasticity).  

\begin{center}
\begin{tabular}{c|c|c|c}
$k$ & $w_k$ & $\nabla f(w_k)$ & $w_{k+1}=w_k-\alpha\nabla f(w_k)$ \\ \hline
0 & $0$ & $-6$ & $0-0.1(-6)=0.6$ \\
1 & $0.6$ & $-4.8$ & $0.6-0.1(-4.8)=1.08$ \\
2 & $1.08$ & $-3.84$ & $1.08-0.1(-3.84)=1.464$ \\
\end{tabular}
\end{center}
Corresponding objective values:
$f(w_0)=9$, $f(w_1)=5.76$, $f(w_2)=3.317\ldots$,
showing a monotone decrease.

%%%%%%%%%%%%%%%%%%%%%%%%%%%%%%%%%%%%%%%%%%%%%%%%%%%%%%%%%%%%
\section*{Assumptions}
\begin{assumption}[Smoothness]\label{ass:smooth}
$f$ is $L$‑smooth, i.e.
\[
\forall x,y\in\R^d,\qquad 
f(y)\le f(x)+\ip{\nabla f(x)}{y-x}+\frac{L}{2}\norm{y-x}^2 .
\tag{2}
\]
\end{assumption}
\begin{assumption}[Lower Boundedness]\label{ass:lower}
There exists $f_{\inf}>-\infty$ such that $f(x)\ge f_{\inf}$ for all $x$.
\end{assumption}
\begin{assumption}[Gradient Growth]\label{ass:gradbound}
The gradient is at most linearly growing:
\[
\forall x\in\R^d,\qquad \norm{\nabla f(x)}\le G\bigl(1+\norm{x}\bigr)
\tag{3}
\]
for some $G\ge0$.
\end{assumption}
\begin{assumption}[Noise Schedule]\label{ass:noise}
The additive noise $\xi_k\sim\mathcal N(0,\sigma_k^2 I_d)$ satisfies
\[
\sigma_k^2 = \frac{c}{(k+1)^{\beta}},\qquad c>0,\;\beta\in (0,1].
\tag{4}
\]
\end{assumption}
\begin{assumption}[Step‑size Schedule]\label{ass:stepsize}
\[
\alpha_k = \frac{\alpha}{(k+1)^{\gamma}},\qquad
\alpha>0,\;\gamma\in\bigl[0,1\bigr].
\tag{5}
\]
\end{assumption}

Throughout the analysis we assume that the sequences
$\{\alpha_k\}$ and $\{\sigma_k\}$ are chosen such that
$\sum_{k=0}^{\infty}\alpha_k=\infty$ and
$\sum_{k=0}^{\infty}\alpha_k\sigma_k^2<\infty$ (the latter will hold
whenever $\gamma+\beta>1$).

%%%%%%%%%%%%%%%%%%%%%%%%%%%%%%%%%%%%%%%%%%%%%%%%%%%%%%%%%%%%
\section*{Main Convergence Theorem}
\begin{theorem}[Expected Gradient Norm Decay]\label{thm:main}
Let Assumptions \ref{ass:smooth}–\ref{ass:stepsize} hold and
choose $\gamma\in[0,1]$, $\beta\in(0,1]$ such that $\gamma+\beta>1$.
Define the averaged iterate
\[
\bar x_T:=\frac{1}{S_T}\sum_{k=0}^{T-1}\alpha_k x_k,
\qquad 
S_T:=\sum_{k=0}^{T-1}\alpha_k .
\]
Then for any $T\ge1$,
\[
\frac{1}{S_T}\sum_{k=0}^{T-1}\alpha_k
\E\!\bigl[\norm{\nabla f(x_k)}^2\bigr]
\;\le\;
\frac{2\bigl(f(x_0)-f_{\inf}\bigr)}{S_T}
\;+\;
\frac{L\alpha^2}{S_T}\sum_{k=0}^{T-1}\frac{1}{(k+1)^{2\gamma}}
\;+\;
\frac{2d\,c}{S_T}\sum_{k=0}^{T-1}\frac{1}{(k+1)^{\gamma+\beta}} .
\tag{6}
\]
In particular, for the concrete choices
$\gamma=\beta=\tfrac12$ (i.e. $\alpha_k=\alpha/\sqrt{k+1}$,
$\sigma_k^2=c/\sqrt{k+1}$) we obtain
\[
\frac{1}{T}\sum_{k=0}^{T-1}
\E\!\bigl[\norm{\nabla f(x_k)}^2\bigr]
\;\le\;
\underbrace{\frac{2\bigl(f(x_0)-f_{\inf}\bigr)}{\alpha\sqrt{T}}}_{\text{bias term}}
\;+\;
\underbrace{\frac{L\alpha}{\sqrt{T}}}_{\text{step‑size term}}
\;+\;
\underbrace{\frac{2d\,c}{\alpha \,T^{1/2}}}_{\text{noise term}} .
\tag{7}
\]
Thus the average squared gradient norm decays as $\mathcal O(T^{-1/2})$.
\end{theorem}

%%%%%%%%%%%%%%%%%%%%%%%%%%%%%%%%%%%%%%%%%%%%%%%%%%%%%%%%%%%%
\section*{Theoretical Properties}
\begin{itemize}
\item \textbf{Stability.}  Because the noise variance decays faster than the
step size (condition $\gamma+\beta>1$), the stochastic perturbation
does not accumulate; the iterates remain bounded in expectation
(see Lemma~\ref{lem:boundedness}).

\item \textbf{Hyper‑parameter Sensitivity.}
  The bound (6) makes explicit how the constants $\alpha$, $c$, $L$ and
  the decay exponents $\gamma,\beta$ trade off bias (first term) versus
  variance (third term).  Larger $\alpha$ accelerates descent but
  inflates the $L\alpha^2$ term; larger $c$ yields a higher steady‑state
  noise floor.

\item \textbf{Comparison to Baselines.}
  - For deterministic gradient descent ($\sigma_k\equiv0$) the third term
    vanishes and we recover the classic $O(1/\sqrt{T})$ rate for
    non‑convex $L$‑smooth objectives.  
  - For Stochastic Gradient Langevin Dynamics with constant temperature
    ($\sigma_k\equiv\sigma>0$) the sum $\sum\alpha_k\sigma_k^2$ diverges,
    preventing convergence to a stationary point; instead the chain
    converges to a stationary distribution.  Our decreasing schedule
    restores pointwise convergence.

\end{itemize}

%%%%%%%%%%%%%%%%%%%%%%%%%%%%%%%%%%%%%%%%%%%%%%%%%%%%%%%%%%%%
\section*{Discussion}
The theorem shows that NGD‑DT enjoys the same $T^{-1/2}$ convergence
rate as classical SGD for non‑convex smooth objectives, while the
added noise can be used for escaping shallow local minima.  By
decreasing the temperature according to (4) the algorithm eventually
behaves like deterministic gradient descent, guaranteeing that the
expected gradient norm vanishes.

%%%%%%%%%%%%%%%%%%%%%%%%%%%%%%%%%%%%%%%%%%%%%%%%%%%%%%%%%%%%
\section*{Preliminaries (Definitions and Lemmas)}
\begin{lemma}[Smoothness Inequality]\label{lem:smooth}
If $f$ is $L$‑smooth then for any $x,y$,
\[
f(y)\le f(x)+\ip{\nabla f(x)}{y-x}+\frac{L}{2}\norm{y-x}^2 .
\tag{8}
\]
\end{lemma}
\begin{proof}
Standard; follows from integrating the Lipschitz condition on
$\nabla f$. \hfill $\square$
\end{proof}

\begin{lemma}[Bounded Second Moment of Noise]\label{lem:noise}
For $\xi_k\sim\mathcal N(0,\sigma_k^2 I_d)$,
\[
\E\!\bigl[\norm{\xi_k}^2\bigr]=d\,\sigma_k^2 .
\tag{9}
\]
\end{lemma}
\begin{proof}
Each coordinate has variance $\sigma_k^2$, sum over $d$ independent
coordinates. \hfill $\square$
\end{proof}

\begin{lemma}[Boundedness of Iterates]\label{lem:boundedness}
Under Assumptions \ref{ass:smooth}–\ref{ass:stepsize} with
$\gamma+\beta>1$ there exists a constant $M>0$ such that
\[
\sup_{k\ge0}\E\!\bigl[\norm{x_k}^2\bigr]\le M .
\tag{10}
\]
\end{lemma}
\begin{proof}
We prove by induction using (1) and Lemma~\ref{lem:smooth}.  From (1)
\[
x_{k+1}=x_k-\alpha_k\nabla f(x_k)+\xi_k .
\]
Taking norms squared, expanding, and using $\E[\xi_k]=0$,
\[
\begin{aligned}
\E\!\bigl[\norm{x_{k+1}}^2\bigr]
&=\E\!\bigl[\norm{x_k-\alpha_k\nabla f(x_k)}^2\bigr]
   +\E\!\bigl[\norm{\xi_k}^2\bigr] \\[0.2em]
&\le \E\!\bigl[\norm{x_k}^2\bigr]
   -2\alpha_k\E\!\bigl[\ip{x_k}{\nabla f(x_k)}\bigr]
   +\alpha_k^2\E\!\bigl[\norm{\nabla f(x_k)}^2\bigr]
   +d\sigma_k^2 .
\end{aligned}
\tag{11}
\]
Using the gradient growth (3) and the fact that $-2\alpha_k\ip{x_k}{\nabla f(x_k)}
\le 2\alpha_k G \norm{x_k}(1+\norm{x_k})$, one obtains a linear
recurrence of the form
\[
\E\!\bigl[\norm{x_{k+1}}^2\bigr]\le (1+C\alpha_k)\E\!\bigl[\norm{x_k}^2\bigr]
      +C'(\alpha_k^2+d\sigma_k^2) .
\tag{12}
\]
Since $\sum_k\alpha_k<\infty$ and $\sum_k\sigma_k^2\alpha_k<\infty$
(the latter holds because $\gamma+\beta>1$), a discrete Grönwall lemma
yields a uniform bound $M$.  \hfill $\square$
\end{proof}

%%%%%%%%%%%%%%%%%%%%%%%%%%%%%%%%%%%%%%%%%%%%%%%%%%%%%%%%%%%%
\section*{Main Proof (Step‑by‑Step Derivation)}
We now prove Theorem~\ref{thm:main}.  All expectations are taken with
respect to the randomness up to time $k$ unless otherwise noted.

\noindent\textbf{[Step 1]} Apply Lemma~\ref{lem:smooth} with
$y=x_{k+1}$, $x=x_k$:
\[
f(x_{k+1})\le f(x_k)+\ip{\nabla f(x_k)}{x_{k+1}-x_k}
            +\frac{L}{2}\norm{x_{k+1}-x_k}^2 .
\tag{13}
\]

\noindent\textbf{[Step 2]} Insert the update rule (1):
\[
x_{k+1}-x_k=-\alpha_k\nabla f(x_k)+\xi_k .
\tag{14}
\]

\noindent\textbf{[Step 3]} Substitute (14) into (13):
\[
\begin{aligned}
f(x_{k+1}) &\le f(x_k)
            +\ip{\nabla f(x_k)}{-\alpha_k\nabla f(x_k)+\xi_k}
            +\frac{L}{2}\norm{-\alpha_k\nabla f(x_k)+\xi_k}^2   \\
           &= f(x_k)-\alpha_k\norm{\nabla f(x_k)}^2
              +\ip{\nabla f(x_k)}{\xi_k}
              +\frac{L}{2}\Bigl(\alpha_k^2\norm{\nabla f(x_k)}^2
                               -2\alpha_k\ip{\nabla f(x_k)}{\xi_k}
                               +\norm{\xi_k}^2\Bigr) .
\end{aligned}
\tag{15}
\]

\noindent\textbf{[Step 4]} Rearrange (15) and collect terms:
\[
\begin{aligned}
f(x_{k+1}) &\le f(x_k)
            -\Bigl(\alpha_k-\tfrac{L}{2}\alpha_k^2\Bigr)\norm{\nabla f(x_k)}^2 \\
           &\quad +\Bigl(1-L\alpha_k\Bigr)\ip{\nabla f(x_k)}{\xi_k}
            +\frac{L}{2}\norm{\xi_k}^2 .
\end{aligned}
\tag{16}
\]

\noindent\textbf{[Step 5]} Take total expectation and use
$\E[\xi_k\,|\,x_k]=0$:
\[
\E\!\bigl[f(x_{k+1})\bigr]\le
\E\!\bigl[f(x_k)\bigr]
-\Bigl(\alpha_k-\tfrac{L}{2}\alpha_k^2\Bigr)
   \E\!\bigl[\norm{\nabla f(x_k)}^2\bigr]
+\frac{L}{2}\E\!\bigl[\norm{\xi_k}^2\bigr] .
\tag{17}
\]

\noindent\textbf{[Step 6]} Apply Lemma~\ref{lem:noise} to replace the
second moment of $\xi_k$:
\[
\E\!\bigl[\norm{\xi_k}^2\bigr]=d\,\sigma_k^2 .
\tag{18}
\]

\noindent\textbf{[Step 7]} Insert (18) into (17) and rearrange:
\[
\Bigl(\alpha_k-\tfrac{L}{2}\alpha_k^2\Bigr)
   \E\!\bigl[\norm{\nabla f(x_k)}^2\bigr]
\le
\E\!\bigl[f(x_k)-f(x_{k+1})\bigr]
+\frac{L d}{2}\sigma_k^2 .
\tag{19}
\]

\noindent\textbf{[Step 8]} Since $0<\alpha_k\le \alpha$, we have
$\alpha_k-\tfrac{L}{2}\alpha_k^2\ge \tfrac12\alpha_k$ provided
$\alpha\le 1/L$ (otherwise replace $\tfrac12$ by a generic constant
$\kappa\in(0,1)$).  For simplicity we assume $\alpha\le 1/L$ and obtain
\[
\frac{\alpha_k}{2}\,
\E\!\bigl[\norm{\nabla f(x_k)}^2\bigr]
\le
\E\!\bigl[f(x_k)-f(x_{k+1})\bigr]
+\frac{L d}{2}\sigma_k^2 .
\tag{20}
\]

\noindent\textbf{[Step 9]} Sum (20) over $k=0,\dots,T-1$:
\[
\frac12\sum_{k=0}^{T-1}\alpha_k
   \E\!\bigl[\norm{\nabla f(x_k)}^2\bigr]
\le
f(x_0)-\E\!\bigl[f(x_T)\bigr]
+\frac{L d}{2}\sum_{k=0}^{T-1}\sigma_k^2 .
\tag{21}
\]

\noindent\textbf{[Step 10]} Drop the non‑negative term
$\E[f(x_T)]\ge f_{\inf}$ and use the explicit schedules
(4)–(5):
\[
\sum_{k=0}^{T-1}\alpha_k
   \E\!\bigl[\norm{\nabla f(x_k)}^2\bigr]
\le
2\bigl(f(x_0)-f_{\inf}\bigr)
+L d\sum_{k=0}^{T-1}\frac{c}{(k+1)^{\beta}} .
\tag{22}
\]

\noindent\textbf{[Step 11]} To obtain the bound in Theorem~\ref{thm:main}
we keep the exact $\alpha_k$ factor inside the noise term, i.e.
multiply (20) by $2/\alpha_k$ before summation:
\[
\begin{aligned}
\sum_{k=0}^{T-1}\alpha_k
   \E\!\bigl[\norm{\nabla f(x_k)}^2\bigr]
&\le
2\bigl(f(x_0)-f_{\inf}\bigr)
+L\alpha^2\sum_{k=0}^{T-1}\frac{1}{(k+1)^{2\gamma}}
+2d\,c\sum_{k=0}^{T-1}\frac{1}{(k+1)^{\gamma+\beta}} .
\end{aligned}
\tag{23}
\]
Dividing (23) by $S_T:=\sum_{k=0}^{T-1}\alpha_k$ yields exactly
the inequality (6).  This completes the proof. \hfill $\square$

%%%%%%%%%%%%%%%%%%%%%%%%%%%%%%%%%%%%%%%%%%%%%%%%%%%%%%%%%%%%
\section*{Rate Derivation (With Constants)}
From standard bounds on harmonic series,
\[
\sum_{k=0}^{T-1}\frac{1}{(k+1)^{p}}
\;\le\;
\begin{cases}
\displaystyle \frac{T^{1-p}}{1-p}, & p\in(0,1),\\[0.4em]
\displaystyle \log T +1, & p=1 .
\end{cases}
\tag{24}
\]
Applying (24) to each term in (6) with $\gamma=\beta=\tfrac12$ gives
\[
S_T = \sum_{k=0}^{T-1}\frac{\alpha}{\sqrt{k+1}}
      \;\ge\;
\alpha\bigl(2\sqrt{T}-2\bigr) = \Theta(\alpha\sqrt{T}) .
\]
Similarly,
\[
\sum_{k=0}^{T-1}\frac{1}{(k+1)^{2\gamma}}=
\sum_{k=0}^{T-1}\frac{1}{k+1}= \log T + O(1),
\qquad
\sum_{k=0}^{T-1}\frac{1}{(k+1)^{\gamma+\beta}}
 =\sum_{k=0}^{T-1}\frac{1}{k+1}= \log T + O(1).
\]
Plugging these estimates into (6) yields (7), i.e. a
$\mathcal O(T^{-1/2})$ decay of the averaged squared gradient norm.

%%%%%%%%%%%%%%%%%%%%%%%%%%%%%%%%%%%%%%%%%%%%%%%%%%%%%%%%%%%%
\section*{Special Cases}
\begin{enumerate}
\item \textbf{Deterministic Gradient Descent.}
  Setting $c=0$ (no noise) reduces (6) to
  \[
  \frac{1}{S_T}\sum_{k=0}^{T-1}\alpha_k
  \E\!\bigl[\norm{\nabla f(x_k)}^2\bigr]
  \le
  \frac{2\bigl(f(x_0)-f_{\inf}\bigr)}{S_T}
  +\frac{L\alpha^2}{S_T}\sum_{k=0}^{T-1}\frac{1}{(k+1)^{2\gamma}} ,
  \]
  which is the classical $O(T^{-1/2})$ bound for non‑convex GD.

\item \textbf{Constant Temperature SGLD.}
  If $\sigma_k\equiv\sigma>0$ (i.e. $\beta=0$) then the series
  $\sum_{k}\alpha_k\sigma_k^2$ diverges; the right‑hand side of (6)
  grows linearly with $T$, and no vanishing gradient guarantee can be
  obtained.  This matches the known behaviour that constant‑temperature
  Langevin dynamics converges to a stationary distribution rather
  than a point.

\item \textbf{Fast Decay of Step‑size.}
  Choosing $\gamma=1$ (step size $\alpha_k=\alpha/(k+1)$) yields
  $S_T = \alpha\log T$, and the bias term in (7) becomes
  $O((\log T)/T)$, i.e. a faster $O(\log T /T)$ decay at the expense of
  a larger $L\alpha^2\sum (k+1)^{-2}$ term, which stays bounded.
\end{enumerate}

%%%%%%%%%%%%%%%%%%%%%%%%%%%%%%%%%%%%%%%%%%%%%%%%%%%%%%%%%%%%
\section*{Conclusion}
We have presented a complete, step‑by‑step convergence analysis for the
noisy gradient descent scheme (1) with decreasing temperature.  Under
standard smoothness assumptions and with step‑size/noise schedules that
satisfy $\gamma+\beta>1$, the method attains the optimal
$\mathcal O(T^{-1/2})$ rate for the expected squared gradient norm in the
non‑convex setting, while retaining the exploratory benefits of
Langevin‑type perturbations.

\end{document}